\begin{abstract}
The geometry of a probabilistic decoder is usually read from its Jacobian.
This excludes black-box decoders and leaves a more basic inverse question
unanswered: do local distances between output distributions determine the
latent surface itself?  We give a positive answer under an explicit convex
polyhedral model.  On a known sphere-type triangulation, local
2-Wasserstein distances define the edge lengths of a piecewise-Euclidean
metric.  The same lengths give exact PL curvature masses by angle defects.
When they are the edges of an unknown strictly convex realization with the
same triangular complex, face congruence and rigidity identify that
realization in $\mathbb R^3$, up to rigid motion.  No decoder derivative or
three-dimensional label is used.

We quantify the metric and local-realization parts of this chain.  Under shape
regularity, additive edge
error $\varepsilon$ perturbs a vertex curvature mass by
$O(d_{\max}\varepsilon/\ell_{\min})$ and its area-normalized curvature by
$O(\varepsilon/h^3)$ on a quasi-uniform mesh of scale $h$.  Relative edge
error controls every intrinsic distance and hence Gromov--Hausdorff and
Gromov--Wasserstein discrepancies.  A rigidity-matrix condition converts edge
error into aligned three-dimensional vertex error locally, while a spectral-gap
condition controls the round-sphere distance-to-coordinate step.  For mesh
families with $s_{X_h}\gtrsim h^\beta$, the resulting Hausdorff tradeoff is
$O(\varepsilon_hh^{-(\beta+1)}+h^2)$.  Our controlled
benchmark hides a unit sphere in a full-rank, noncommuting Gaussian covariance
path: decoder means collapse it to a disk, while local Wasserstein distances recover its full
metric.  A round-sphere sweep separates stable global shape from
noise-sensitive pointwise curvature.  A second edge-only realization recovers
the convex hull of samples from a nonround ellipsoid without a spherical final
constraint.
\end{abstract}

\section{Introduction}

A flat map does not carry its own scale.  Equal displacements near the equator
and near a pole can represent different travel distances.  If sufficiently
many local travel times are known, however, the missing scale, curvature, and
eventually the globe can be recovered.  This paper studies the analogous
problem for a generative decoder.  A latent code is a map coordinate, a decoded
probability distribution is the observation at that coordinate, and local
optimal-transport costs are the measured travel times.

Prior work commonly starts from a differentiable decoder and pulls an ambient
metric back with its Jacobian \cite{arvanitidis2018,chen2018,roy2022}.
That construction is useful once the decoder is known, but it does not say how
much geometry is determined by finitely many distance queries.  It also makes
the visual output ambiguous: a metric tensor or a curvature heat map is not
yet a reconstructed surface.

We instead pose an explicit inverse problem.

\begin{quote}
Given a triangulated latent surface and noisy $W_2$ distances only on its edges,
how accurately can one recover intrinsic curvature and the hidden 3D surface?
\end{quote}

The input and output are deliberately finite.  Let
$\mathcal T=(V,E,F)$ be a closed triangular complex homeomorphic to $S^2$.
Let $(\mathcal Y,d_{\mathcal Y})$ be a common complete separable ground space,
and let $W_2$ use the squared ground cost $d_{\mathcal Y}^2$.  At vertex $i$
the decoder returns $\mu_i\in\mathcal P_2(\mathcal Y)$, but the estimator
receives only
\begin{equation}
 \mathcal O=
 \left(\mathcal T,
 \{\widehat\ell_{ij}:(i,j)\in E\}\right),
 \qquad
 \widehat\ell_{ij}\simeq W_2(\mu_i,\mu_j).
 \label{eq:observations}
\end{equation}
Ground-truth points in $\mathbb R^3$ are withheld until evaluation.  The edge
lengths make every face a Euclidean triangle and therefore define a complete
piecewise-Euclidean (PL) metric.  All quantities in the reconstruction---face
angles, curvature, area, and intrinsic distance---are functions of
\eqref{eq:observations} alone.
All finite identifiability and stability results below use only $E$.  The
round-sphere visualization later augments it with a three-hop local query set
$E_{\mathrm q}\supset E$, still of linear size, to approximate global travel
times; those additional queries are not used by the same-skeleton theorem.

For extrinsic recovery we make a separate assumption instead of defining an
ambient target after seeing the observations.

\begin{definition}[Same-skeleton convex observation model]
\label{def:same-skeleton-model}
There is an unknown nondegenerate convex simplicial polyhedron
$X_*=(x_i)_{i\in V}\subset\mathbb R^3$ whose boundary complex is exactly
$\mathcal T$, and
\begin{equation}
 \ell_{ij}=W_2(\mu_i,\mu_j)=\norm{x_i-x_j}
 \qquad ((i,j)\in E).
 \label{eq:convex-observation-model}
\end{equation}
Every vertex is extreme and every interior dihedral angle is strictly below
$\pi$; we call this a strictly convex realization.  The measurements perturb
the displayed lengths.  The complex and convex model class are known;
$X_*$ is not.
\end{definition}

Without Definition~\ref{def:same-skeleton-model}, the edge observations still
define a valid PL proxy $M_\ell$, but need not recover a pre-existing ambient
surface.  An arbitrary endpoint $W_2$ is an ambient chord, not an intrinsic
geodesic.  Proposition~\ref{prop:hidden-sphere} verifies
\eqref{eq:convex-observation-model} exactly for our benchmark.

\paragraph{Contributions.}
\begin{enumerate}
  \item We formulate \emph{Wasserstein travel-time reconstruction}: local
  decoder-distribution distances define a PL latent metric without evaluating
  a decoder Jacobian.
  \item We connect three identifiable levels under a stated same-skeleton
  convex model.  Edge lengths determine angle-defect curvature masses;
  congruent triangular faces determine the exact convex realization; and a
  nondegenerate rigidity matrix makes that realization locally stable.
  \item We give explicit perturbation bounds from edge noise to curvature,
  every intrinsic distance, GH/GW discrepancy, and aligned vertex error.  The
  bounds expose the noise sensitivity of curvature: at mesh scale $h$,
  additive length noise is amplified as $\varepsilon/h^3$ after area
  normalization.  A conditional refinement corollary turns rigidity scaling
  into an explicit resolution--noise requirement.
  \item We introduce a ground-truth benchmark in which the decoder mean shows
  only a disk while a full-rank noncommuting covariance path stores the missing
  height.  Gaussian $W_2$ distances recover the exact edge lengths of an inscribed sphere, so the
  probabilistic observation model is exact and auditable.  An edge-only
  continuation experiment then verifies the same-skeleton step on an
  ellipsoidal convex polyhedron.
\end{enumerate}

\paragraph{Scope of the OT claim.}
Once edge lengths are supplied, our curvature, metric-stability, and rigidity
results are metric-input statements and are not specific to optimal transport.
The controlled model does use the Bures part of Gaussian $W_2$: its full-rank
covariances do not commute, and their raw Frobenius distance has the
position-dependent distortion in
\eqref{eq:covariance-frobenius-distortion}.  It is nevertheless a designed
Wasserstein geodesic, not a proof that $W_2$ is uniquely necessary.  Our claimed
contribution is the end-to-end chain from a queryable decoder-distribution
metric to PL curvature and a convex inverse reconstruction, not a new general
graph-realization method.

The paper does not infer topology, assert that a general intrinsic metric has a
unique nonconvex embedding, or claim a true curvature for MNIST.  Sphere
topology and convexity are assumptions, while the smooth sphere is a controlled
evaluation target.

\section{Related Work}

\paragraph{Latent geometry.}
Pullback metrics for deterministic and stochastic generative models use
decoder derivatives to define latent lengths and geodesics
\cite{arvanitidis2018,chen2018}.  Optimal Latent Transport pulls a Wasserstein
metric back from decoder distributions \cite{roy2022}.  Syrota et
al.~prove that distances, angles, and volumes of deep latent-variable models
can be identifiable even when coordinates are not \cite{syrota2025}.  These
works establish or exploit geometry for a known differentiable model.  Our
target is different: reconstruction accuracy from sparse, derivative-free
Wasserstein distance observations.

\paragraph{Distances to manifolds.}
Qiu and Dai estimate a smooth Riemannian metric and its derivatives from noisy
intrinsic-distance responses, including a projected sphere, taxi travel time,
and MNIST examples \cite{qiu2023}.  Fefferman et al.~reconstruct an abstract
bounded-geometry manifold from noisy intrinsic distances with finite-sample
guarantees \cite{fefferman2020}.  Consequently, distance-to-metric recovery by
itself is not our novelty.  We specialize the observations to local decoder
$W_2$ queries and continue from an intrinsic metric to discrete curvature and
a convex extrinsic realization.  Isomap also chains local distances, but its
goal is a Euclidean coordinate representation rather than curvature or convex
surface identification \cite{tenenbaum2000}.
Hamm et al.~recover the intrinsic metric and tangent spaces of a sampled
submanifold of Wasserstein space in an asymptotic GW sense
\cite{hamm2025}.  We instead take a finite PL target and continue to its
curvature masses and identifiable convex realization.

\paragraph{Euclidean distance geometry.}
Recovering Euclidean coordinates from incomplete distances is the classical
distance-geometry problem; its theory includes graph rigidity, matrix
completion, and sensor localization \cite{liberti2014}.  Our edge-only
realization is an instance of that problem, not a new graph-realization method.
The distinction is the preceding observation-and-guarantee chain: local
distances between decoder distributions define a PL metric and curvature
before a stated convex model permits an extrinsic realization.

\paragraph{Polyhedral geometry and shape comparison.}
Alexandrov's theorem characterizes convex polyhedra by their intrinsic PL
metrics; Bobenko and Izmestiev give a constructive proof
\cite{alexandrov2005,bobenko2008}.  Angle defect is the intrinsic curvature
measure of a PL surface.  Its approximation of smooth Gaussian curvature
requires mesh assumptions and is not pointwise consistent on arbitrary
triangulations \cite{gawlik2020,hartig2021}.  We therefore state the finite-
noise theorem first for the exact PL target and report smooth-sphere error
separately.  Gromov--Wasserstein distance compares metric-measure spaces up to
measure-preserving isometry \cite{memoli2011}; a known vertex correspondence
also supplies a simpler upper bound through its distance distortion.

\section{Wasserstein Edge Geometry}
\label{sec:problem}

Let every face $f=(i,j,k)\in F$ satisfy the strict triangle inequalities for
the positive lengths $\ell_{ij},\ell_{jk},\ell_{ki}$.  Gluing the resulting
Euclidean triangles along equal edges defines a PL metric space
$M_\ell=(\mathcal T,d_\ell)$.  This object is intrinsic: no coordinates in an
ambient Euclidean space are part of its definition.

For a face incident to $i$, the angle at $i$ is computable by the cosine law,
\begin{equation}
 \theta_i^{jk}(\ell)
 =\arccos\!\left(
 \frac{\ell_{ij}^2+\ell_{ik}^2-\ell_{jk}^2}
 {2\ell_{ij}\ell_{ik}}
 \right).
 \label{eq:cosine-angle}
\end{equation}
The exact PL curvature mass at an interior vertex and its
barycentric-area density are
\begin{align}
 \kappa_i(\ell)
 &=2\pi-\sum_{f\ni i}\theta_{i,f}(\ell),
 \label{eq:angle-defect}\\
 A_i(\ell)
 &=\frac13\sum_{f\ni i}A_f(\ell),
 &K_i(\ell)&=\frac{\kappa_i(\ell)}{A_i(\ell)}.
 \label{eq:discrete-density}
\end{align}
Here $\kappa_i$ is an integrated curvature mass, measured in radians; it is
not a pointwise Gaussian curvature.  The normalized $K_i$ is useful for
comparison with a smooth surface, but depends additionally on the chosen
vertex area and need not be pointwise consistent on an arbitrary mesh.  For
every closed sphere triangulation,
\begin{equation}
 \sum_{i\in V}\kappa_i=4\pi
 \label{eq:discrete-gauss-bonnet}
\end{equation}
by the Euclidean triangle angle sum and Euler's formula.

\subsection{A sphere hidden in Gaussian uncertainty}

The following construction is the controlled instance used in the experiment.
It also makes clear why a mean-only distance is insufficient.

\begin{proposition}[Noncommuting hidden-height Wasserstein identity]
\label{prop:hidden-sphere}
Fix $R>0$ and $k\in(0,1)$, and let
\begin{equation}
 B=\begin{pmatrix}1&0\\0&2\end{pmatrix},\qquad
 H=k\begin{pmatrix}0&1\\1&0\end{pmatrix},\qquad
 \lambda=\frac{4R^2}{3k^2}.
 \label{eq:hidden-decoder}
\end{equation}
For $x=(x_1,x_2,x_3)$ with $|x_3|\le R$, set
\begin{equation}
 \begin{aligned}
  t_x&=\frac{x_3+R}{2R}, &m_x&=(x_1,x_2),\\
  C_x&=\lambda(I_2+t_xH)B(I_2+t_xH).
 \end{aligned}
 \label{eq:hidden-covariance-path}
\end{equation}
and $\mu_x=\mathcal N(m_x,C_x)$.  Then every $C_x$ is positive definite and,
for every $x,y$ in this slab, the Gaussian Bures covariance distance
$d_{\mathrm B}$ satisfies
\begin{equation}
 \begin{aligned}
 W_2^2(\mu_x,\mu_y)
 &=\norm{m_x-m_y}^2+d_{\mathrm B}^2(C_x,C_y)\\
 &=\norm{x-y}^2.
 \end{aligned}
 \label{eq:hidden-identity}
\end{equation}
If $x_3\ne y_3$, then $C_xC_y\ne C_yC_x$.  On $S_R^2$, the mean map is a
disk that identifies northern and southern points with the same
$(x_1,x_2)$; the Wasserstein observation distinguishes them.
\end{proposition}

\begin{proof}
Write $A_t=I_2+tH$ and $L_t=\sqrt\lambda A_tB^{1/2}$, so
$C_t=L_tL_t^\trans$.  Since $A_s$ and $A_t$ are commuting positive definite
polynomials in $H$, $L_s^\trans L_t$ is positive definite.  The orthogonal
Procrustes formula for the Gaussian Bures distance therefore selects the
identity alignment and gives \cite{givensshortt1984,takatsu2011}
\begin{equation}
 \begin{aligned}
 d_{\mathrm B}^2(C_s,C_t)
 &=\norm{L_s-L_t}_F^2\\
 &=3\lambda k^2(s-t)^2
 =4R^2(s-t)^2.
 \end{aligned}
 \label{eq:bures-height-identity}
\end{equation}
Substituting $s=t_x$ and $t=t_y$ gives
$d_{\mathrm B}(C_x,C_y)=|x_3-y_3|$, proving \eqref{eq:hidden-identity}.
Positive definiteness follows from $t_x\in[0,1]$ and $k<1$.
Direct multiplication gives
\[
 C_sC_t-C_tC_s
 =3\lambda^2k(s-t)(1+k^2st)
 \begin{pmatrix}0&1\\-1&0\end{pmatrix},
\]
which is nonzero for $s\ne t$.  In contrast, raw covariance coordinates obey
\begin{equation}
 \norm{C_s-C_t}_F
 =\lambda k|s-t|\sqrt{18+5k^2(s+t)^2},
 \label{eq:covariance-frobenius-distortion}
\end{equation}
so no single rescaling of their Frobenius distance recovers the height metric.
\end{proof}

If two sphere points have great-circle distance $d\in[0,\pi R]$, their
Wasserstein distance in \eqref{eq:hidden-identity} is the chord
\begin{equation}
 W_2(\mu_x,\mu_y)=2R\sin\!\left(\frac{d}{2R}\right),
 \qquad
 0\le d-W_2\le \frac{d^3}{24R^2}.
 \label{eq:chord-arc}
\end{equation}
Thus a single remote $W_2$ query takes a shortcut through the sphere.  A chain
of local queries removes that shortcut as the mesh is refined.  This is why
our observation graph contains local edges rather than all endpoint pairs.

\section{What Can Be Recovered?}
\label{sec:theory}

We separate exact identifiability from stability under noise.  The first
result is classical, but it supplies the missing condition behind the phrase
``recover the surface.''

\begin{theorem}[Exact same-skeleton identification]
\label{thm:convex-identification}
Under Definition~\ref{def:same-skeleton-model}, exact local Wasserstein edge
lengths determine $X_*$ up to an element of the Euclidean group.
\end{theorem}

\begin{proof}
Each corresponding triangular face is congruent by the side--side--side
criterion.  Cauchy's rigidity theorem then makes two convex polyhedra with
congruent corresponding faces congruent \cite{alexandrov2005,pak2006}.
\end{proof}

This theorem is conditional on existence of the same-complex convex target;
it does not infer convexity or combinatorics.  Alexandrov's theorem gives a
different, purely intrinsic realization result for suitable positive-curvature
PL metrics \cite{bobenko2008}, but its natural polyhedral edges need not equal
the prescribed input graph.  We therefore do not compose Alexandrov's theorem
with the edge-rigidity bound below.

\subsection{Edge noise to curvature noise}

Let $d_i$ be the number of faces incident to vertex $i$ and set
$d_{\max}=\max_i d_i$.  We use the following shape-regular class.  Both the
true and perturbed triangles have sides in
$[\ell_{\min}/2,2\ell_{\max}]$ and angles at least $\alpha_0>0$; set
$\rho=\ell_{\max}/\ell_{\min}$.  These conditions exclude triangles that are
arbitrarily close to degenerate.

\begin{theorem}[Stability of discrete curvature]
\label{thm:curvature-noise}
Suppose
\begin{equation}
 \max_{e\in E}\abs{\widehat\ell_e-\ell_e}
 \le\varepsilon.
 \label{eq:additive-noise}
\end{equation}
Under the shape-regular assumptions above, constants
$C_\theta=C_\theta(\rho,\alpha_0)$ and
$C_A=C_A(\rho,\alpha_0)$ satisfy
\begin{align}
 \abs{\widehat\kappa_i-\kappa_i}
 &\le C_\theta d_i
       \frac{\varepsilon}{\ell_{\min}},
 \label{eq:defect-stability}\\
 \abs{\widehat A_i-A_i}
 &\le C_A d_i\ell_{\max}\varepsilon.
 \label{eq:area-stability}
\end{align}
If the right side of \eqref{eq:area-stability} is at most $A_i/2$, then
\begin{equation}
 \abs{\widehat K_i-K_i}
 \le
 \frac{2C_\theta d_i\varepsilon}
      {A_i\ell_{\min}}
 +\frac{2\abs{\kappa_i}C_A d_i\ell_{\max}\varepsilon}
      {A_i^2}.
 \label{eq:density-stability}
\end{equation}
In particular, positive defect masses are retained whenever
\begin{equation}
 C_\theta d_{\max}\frac{\varepsilon}{\ell_{\min}}
 <\kappa_{\min}:=\min_i\kappa_i.
 \label{eq:convexity-noise}
\end{equation}
\end{theorem}

\begin{proof}
On a compact shape-regular set of side triples, differentiating the cosine law
\eqref{eq:cosine-angle} bounds the $\ell_1$ norm of every angle gradient by
$C_\theta/\ell_{\min}$; the factor $1/\sin\theta$ is controlled by
$\alpha_0$.  The mean-value theorem and summation over the incident faces give
\eqref{eq:defect-stability}.  Heron's formula is smooth on the same compact
set, with gradient bounded by $C_A\ell_{\max}$, which proves
\eqref{eq:area-stability}.  Finally,
\[
 \frac{\widehat\kappa_i}{\widehat A_i}
 -\frac{\kappa_i}{A_i}
 =\frac{\widehat\kappa_i-\kappa_i}{\widehat A_i}
 +\kappa_i\frac{A_i-\widehat A_i}{A_i\widehat A_i},
\]
and $\widehat A_i\ge A_i/2$ yields \eqref{eq:density-stability}.
Condition \eqref{eq:convexity-noise} makes every perturbed defect positive.
\end{proof}

On a quasi-uniform mesh with edge scale $h$, bounded degree, and
$A_i\asymp h^2$, Theorem~\ref{thm:curvature-noise} reads
\begin{equation}
 \abs{\widehat\kappa_i-\kappa_i}
 =O(\varepsilon/h),
 \qquad
 \abs{\widehat K_i-K_i}=O(\varepsilon/h^3).
 \label{eq:noise-amplification}
\end{equation}
The first quantity is the exact PL curvature measure.  The second divides by a
shrinking area and is necessarily more noise-sensitive.  Equation
\eqref{eq:noise-amplification} is a perturbation result, not a claim of
pointwise convergence to smooth curvature on arbitrary meshes.  If the edge
error is relative, $\varepsilon\asymp\eta h$, then the two orders become
$O(\eta)$ and $O(\eta/h^2)$.  If, in addition, the mesh family is
curvature-consistent in the sense that $\kappa_i\asymp h^2$, retaining every
defect sign is guaranteed by $\eta=O(h^2)$ with a sufficiently small constant.
Thus
fixed relative noise can preserve global distances while destroying
pointwise curvature under refinement; the experiment reports both effects.

\subsection{Intrinsic and extrinsic error}

Relative length noise yields a scale-free guarantee.  Each pair of
corresponding true and perturbed faces has a canonical affine map in barycentric
coordinates.  Shape regularity makes its singular values stable.

\begin{theorem}[Distance and metric-space stability]
\label{thm:distance-stability}
Assume shape regularity and
\begin{equation}
 \max_{e\in E}
 \abs{\widehat\ell_e/\ell_e-1}\le\eta
 \label{eq:relative-noise}
\end{equation}
for sufficiently small $\eta$.  There is
$C_g=C_g(\rho,\alpha_0)$ such that the facewise affine correspondence
$\Phi:M_\ell\to M_{\widehat\ell}$ obeys
\begin{equation}
 (1-C_g\eta)d_\ell(p,q)
 \le d_{\widehat\ell}(\Phi p,\Phi q)
 \le(1+C_g\eta)d_\ell(p,q).
 \label{eq:all-distance-stability}
\end{equation}
Consequently,
\begin{equation}
 d_{\mathrm{GH}}(M_\ell,M_{\widehat\ell})
 \le\frac{C_g\eta}{2}\diam(M_\ell).
 \label{eq:gh-bound}
\end{equation}
For any probability measure $\nu$ on $M_\ell$, including fixed uniform
vertex weights, the coupling induced by $p\mapsto(p,\Phi p)$ gives
\begin{equation}
 GW_2\!\left((M_\ell,d_\ell,\nu),
 (M_{\widehat\ell},d_{\widehat\ell},\Phi_\#\nu)\right)
 \le C_g\eta\diam(M_\ell).
 \label{eq:gw-bound}
\end{equation}
\end{theorem}

\begin{proof}[Proof sketch]
The Gram matrix of a Euclidean triangle is a smooth function of its three
squared side lengths.  On the compact shape-regular class,
\eqref{eq:relative-noise} therefore bounds the singular values of the affine
map by $1\pm C_g\eta$.  Mapping each rectifiable path face by face, integrating
the length bounds, and taking infima proves
\eqref{eq:all-distance-stability}.  The graph of $\Phi$ is a correspondence
with distortion at most $C_g\eta\diam(M_\ell)$, giving
\eqref{eq:gh-bound}.  Using the induced measure coupling in the definition of
$GW_2$ gives the final claim \cite{memoli2011}.
\end{proof}

Equation~\eqref{eq:gw-bound} deliberately uses the same abstract probability
measure on both metrics.  The two metrics' separately normalized area
measures are not exact pushforwards of one another; comparing those natural
measures requires an additional density-perturbation term.

The round-sphere experiment has one further map: approximate global distances
are converted to three-dimensional coordinates by a spectral factorization.
The following proposition isolates the stability of that final step.

\begin{proposition}[Local stability of spherical spectral reconstruction]
\label{prop:spherical-spectral-stability}
Let $X\in\R^{n\times3}$ have rows $x_i\in S_R^2$, let $G=XX^\trans$, and
assume $\gamma:=\lambda_3(G)>0$, so $\lambda_4(G)=0$.  Define
\begin{equation}
 D_{ij}=R\arccos\!\left(\frac{x_i^\trans x_j}{R^2}\right).
 \label{eq:true-spherical-distance-matrix}
\end{equation}
For a symmetric $\widehat D$ and $\widehat R>0$, form
$\widehat G_{ij}=\widehat R^2\cos(\widehat D_{ij}/\widehat R)$, retain its
three largest eigenvalues, factor the resulting positive semidefinite matrix
as $\widehat Y\widehat Y^\trans$, and set
$\widehat x_i=\widehat R\widehat y_i/\norm{\widehat y_i}$ row by row.  There
is a neighbourhood of $(D,R)$ and a finite
$C_{\mathrm{sph}}=C_{\mathrm{sph}}(X)$ such that the retained eigenvalues are positive,
all rows are nonzero, and
\begin{equation}
 \min_{Q\in O(3)}\norm{\widehat X-XQ}_F
 \le C_{\mathrm{sph}}\left(
 \norm{\widehat D-D}_F+\sqrt n\abs{\widehat R-R}
 \right).
 \label{eq:spherical-spectral-stability}
\end{equation}
No separation among the leading three eigenvalues is required.
\end{proposition}

\begin{proof}[Proof sketch]
For $B(D,r)_{ij}=r^2\cos(D_{ij}/r)$, the scalar derivatives are
$-r\sin(d/r)$ in $d$ and $2r\cos(d/r)+d\sin(d/r)$ in $r$.  They are bounded
near $(D,R)$, so $B$ is locally Lipschitz in the norms appearing in
\eqref{eq:spherical-spectral-stability}.  Weyl's inequality and
$\gamma>0$ separate the leading three-dimensional eigenspace from the
remaining spectrum.  The projector onto this isolated eigenvalue cluster is
a Riesz spectral projector and varies smoothly while the gap stays open;
rank-three positive spectral truncation is therefore locally Lipschitz.
For completeness, align two full-column-rank factors by orthogonal
Procrustes.  On the resulting horizontal slice, the derivative of
$Y\mapsto YY^\trans$ at $X$ is
$E\mapsto XE^\trans+EX^\trans$.  In a singular-vector basis for $X$, a direct
block calculation bounds its gain below by
$\sqrt{2}\,\sigma_3(X)=\sqrt{2\gamma}$; the quadratic remainder $EE^\trans$
can be absorbed locally.  Hence Gram-factor perturbation gives
\[
 \min_{Q\in O(3)}\norm{\widehat Y-XQ}_F
 \le \frac{C_0}{\sqrt\gamma}\norm{\widehat G-G}_F
\]
locally, including when eigenvalues repeat inside the leading subspace.
Finally, for $\norm x=R$ and $y\ne0$,
$\|\widehat R y/\|y\|-x\|\le|\widehat R-R|+2\|y-x\|$.
Applying this row by row proves the result.
\end{proof}

This proposition begins after $\widehat D$ and $\widehat R$ have been
constructed.  It does not by itself guarantee the graph shortest paths,
diameter calibration, or radius estimator used in the controlled benchmark.

Theorem~\ref{thm:distance-stability} is intrinsic.  To measure reconstructed
3D vertices, let $X\in\mathbb R^{|V|\times3}$ be a strictly convex triangulated
realization and
$\mathcal L(X)=(\norm{X_i-X_j})_{(i,j)\in E}$ its edge-length map.  Its
Jacobian is the rigidity matrix $R_X$.  Write $\mathcal Z_X=\ker R_X$ for the
six-dimensional space of infinitesimal Euclidean motions and define, using
the Frobenius norm on coordinates,
\begin{equation}
 s_X=\sigma_{\min}\!\left(R_X\big|_{\mathcal Z_X^\perp}\right).
 \label{eq:rigidity-singular-value}
\end{equation}

\begin{proposition}[Local same-skeleton 3D stability]
\label{prop:rigidity}
For a strictly convex triangulated realization, $s_X>0$.  Every length vector
$\widehat\ell$ sufficiently close to $\mathcal L(X)$ has a unique nearby
strictly convex realization $\widehat X$ with complex $\mathcal T$, modulo
Euclidean motions, and
\begin{equation}
 \min_{Q\in O(3),\,b\in\R^3}
 \norm{\widehat X-(XQ+\mathbf 1b^{\trans})}_{F}
 \le
 \frac{2}{s_X}
 \norm{\widehat\ell-\mathcal L(X)}_2.
 \label{eq:rigidity-bound}
\end{equation}
The factor $2$ can be replaced by $1+o(1)$ as the edge error tends to zero.
\end{proposition}

\begin{proof}[Proof sketch]
Convex triangulated polyhedra are infinitesimally rigid \cite{pak2006}.
For a sphere triangulation, $|E|=3|V|-6$.  Taking the affine slice through $X$
tangent to $\mathcal Z_X^\perp$ therefore makes the derivative of $\mathcal L$
square and nonsingular, with least singular value $s_X$.  The inverse-function
theorem gives the local realization and uniqueness.  Shrinking the
neighbourhood makes the inverse derivative norm at most $2/s_X$ and preserves
strict convexity.  Minimization over rigid motions can only reduce the slice
displacement, proving \eqref{eq:rigidity-bound}.  The result is local:
it does not assert that an unconstrained optimizer will find this branch.
\end{proof}

Together, Theorem~\ref{thm:convex-identification} and
Proposition~\ref{prop:rigidity} distinguish exact global identification from
local numerical conditioning.  The constant $s_X$ is realization-specific;
we do not assume that it stays bounded away from zero under refinement.

\subsection{The smooth sphere is a separate approximation target}

Exact edge recovery identifies a finite polyhedron, not the smooth sphere at
finite resolution.  The remaining error can be bounded independently of the
observation noise.

\begin{proposition}[Inscribed-sphere approximation]
\label{prop:inscribed-sphere}
Let the vertices $x_i$ lie on the radius-$R$ sphere, suppose their convex hull
$P_h$ contains the origin, and let $\mathcal T$ be its simplicial boundary.
If
\begin{equation}
 h=\max_{(i,j)\in E}d_{S_R^2}(x_i,x_j)<\frac{\pi R}{2},
 \label{eq:spherical-mesh-size}
\end{equation}
then exact observations from Proposition~\ref{prop:hidden-sphere} recover the
intrinsic metric of $\partial P_h$.  Moreover,
\begin{equation}
 d_{\mathrm H}(\partial P_h,S_R^2)
 \le R\bigl(1-\cos(h/R)\bigr)
 \le\frac{h^2}{2R}.
 \label{eq:inscribed-sphere-bound}
\end{equation}
If noisy lengths remain in the local neighbourhood of
Proposition~\ref{prop:rigidity}, write $\widehat P$ for the convex polyhedron
with vertices $\widehat X$.  Then, after rigid alignment,
\begin{equation}
 d_{\mathrm H}(\partial\widehat P,S_R^2)
 \le
 \frac{2}{s_X}\norm{\widehat\ell-\ell}_2
 +R\bigl(1-\cos(h/R)\bigr)
 \label{eq:noisy-sphere-decomposition}
\end{equation}
All Hausdorff distances in this proposition use the ambient Euclidean metric.
\end{proposition}

\begin{proof}
Equation~\eqref{eq:hidden-identity} makes every reconstructed face congruent
to the corresponding face of $P_h$; Theorem~\ref{thm:convex-identification}
identifies the same-complex convex realization.  For a face containing
$x_i$, set $u=x_i/R$.  Every vertex $x_j$ of that face satisfies
$u^\trans x_j\ge R\cos(h/R)$; hence every convex combination $q$ in the face
satisfies
\[
 \norm q\ge u^\trans q\ge R\cos(h/R).
\]
Also $\norm q\le R$ because $q$ belongs to the convex hull.  Thus radial
projection moves $q$ by at most $R(1-\cos(h/R))$.  Conversely, because the
origin lies in $P_h$, every ray from the origin meets $\partial P_h$; applying
the same bound to that intersection proves the reverse directed Hausdorff
bound.  The elementary inequality
$1-\cos t\le t^2/2$ proves \eqref{eq:inscribed-sphere-bound}.
The triangle inequality with Proposition~\ref{prop:rigidity} gives
\eqref{eq:noisy-sphere-decomposition}.
\end{proof}

In particular, the uniform error assumption
$\max_e|\widehat\ell_e-\ell_e|\le\varepsilon$ gives
\begin{equation}
 d_{\mathrm H}(\partial\widehat P,S_R^2)
 \le
 \frac{2\sqrt{|E|}}{s_X}\,\varepsilon
 +\frac{h^2}{2R}.
 \label{eq:end-to-end-sphere}
\end{equation}
Equation~\eqref{eq:end-to-end-sphere} is the paper's simplest answer to the
question of how close the recovered manifold is: one term measures local
Wasserstein observation error and the other finite-mesh approximation error.

\begin{corollary}[Conditional resolution--noise tradeoff]
\label{cor:conditional-refinement}
Consider a quasi-uniform family on a fixed sphere for which
$|E_h|\le C_E h^{-2}$ and $s_{X_h}\ge c_s h^\beta$ with $\beta\ge0$.
Assume at every resolution that the reconstruction stays in the local branch
of Proposition~\ref{prop:rigidity}, and denote its convex polyhedron by
$\widehat P_h$.  Then
\begin{equation}
 d_{\mathrm H}(\partial\widehat P_h,S_R^2)
 \le \frac{2\sqrt{C_E}}{c_s}\,
 \varepsilon_hh^{-(\beta+1)}+\frac{h^2}{2R}.
 \label{eq:conditional-refinement-rate}
\end{equation}
Consequently, $\varepsilon_h=o(h^{\beta+1})$ suffices for consistency, and
$\varepsilon_h=O(h^{\beta+3})$ retains the $O(h^2)$ discretization rate.  To
target Hausdorff error $\delta$, it is sufficient to take
$h\asymp\delta^{1/2}$ and
$\varepsilon_h=O(\delta^{(\beta+3)/2})$, using $O(\delta^{-1})$ edges.
Equivalently, balancing the terms for a given small $\varepsilon$ gives error
$O(\varepsilon^{2/(\beta+3)})$.
\end{corollary}

\begin{proof}
Substitute $\sqrt{|E_h|}\le\sqrt{C_E}h^{-1}$ and
$s_{X_h}^{-1}\le c_s^{-1}h^{-\beta}$ into
\eqref{eq:end-to-end-sphere}; the remaining claims follow by comparing or
balancing the two terms.
\end{proof}

The two terms in \eqref{eq:noisy-sphere-decomposition} have different
meanings.  The first is observation and reconstruction error for the finite PL
target; the second is discretization error between that target and the smooth
sphere.  We report them separately rather than presenting a zero-noise
polyhedron as the exact smooth sphere.  The rigidity constant is local to a
fixed polyhedron and is not claimed to remain uniform under refinement.
